\chapter{NoSQL Databases}

This chapter provides a more in-depth description of the NoSQL database types introduced in the first chapter, focusing on their internal mechanisms, design patterns, and practical trade-offs.

\section{Key-value Databases}

A key-value database is the simplest form of NoSQL database, storing data as a collection of key-value pairs.

\begin{definition}{Key}
A \keyword{key} is a unique identifier used to look up a value in the database. Keys should have a logical structure, follow a consistent naming convention, and adopt a standard delimiter (e.g., \code{user:1001:profile}, \code{session:abc123}).
\end{definition}

\subsection{Storage and Distribution}

Sets of key-value pairs are assigned to different nodes in a cluster. \keyword{Hashing} is a common method of assigning key-value pairs to nodes: a hash function is applied to the key, and the result determines which node stores the pair. This enables horizontal scaling and efficient distributed lookups.

\begin{note}{Lookup Efficiency}
    A naive key lookup is no more efficient than a linear search. However, more efficient solutions are possible through the use of \keyword{secondary indexes}, which allow querying by value rather than only by key.
\end{note}

\subsection{Design Patterns}

Several design patterns are commonly used when modeling data with key-value databases:

\begin{itemize}
    \item \textbf{TTL Keys:} Keys can have a Time-To-Live, after which they are automatically expired and deleted. This is useful for session data, caches, and temporary tokens.
    \item \textbf{Emulating Tables:} By using compound keys with a shared prefix (e.g., \code{user:1001:name}, \code{user:1001:email}), one can emulate the structure of a relational table.
    \item \textbf{Aggregates:} Related data that is frequently accessed together can be stored as a single value (e.g., serialized JSON). This avoids multiple round-trips to the database.
    \item \textbf{Atomic Aggregates:} When multiple values form an \keyword{atomic aggregate}, they must be updated together or not at all. Storing them under a single key guarantees atomicity, as operations on a single key are atomic in most key-value stores.
    \item \textbf{Indexes:} To enable querying by non-key attributes, secondary indexes (maintained manually or provided by the database) can be built, mapping attribute values to sets of primary keys.
\end{itemize}

\begin{example}{}
We want to represent a key-value database for vehicle rentals.
Each booking has the following attributes: id, vehicle type (car or van), booking start date, booking end date, customer name, customer id.

The critical operations (to be executed most frequently) are:
\begin{enumerate}
    \item Retrieve all bookings for a given start date.
    \item Retrieve the booking given the booking ID.
    \item Update the end date of a booking given the ID.
    \item Delete a booking given the ID.
\end{enumerate}

\vspace{2mm}
\textbf{Option 1: Simulating a table}.

\begin{verbatim}
rent:id:type = <"Car" | "Van">
rent:id:s_date = <date>
rent:id:e_date = <date>
rent:id:cust = { name: <name>, id: <id> }
\end{verbatim}

The first operation is very inefficient: there is no index on \texttt{s\_date}, so a full scan over all bookings would be required.
The second operation is somewhat slow but acceptable: using the id, all attributes can be accessed, but not atomically.
The third operation is instantaneous.
The fourth operation is somewhat slow for the same reasons as the second.

\vspace{2mm}
\textbf{Option 2}.

\begin{verbatim}
rent:id = { 
    type: <type>, 
    s_date: <data>, 
    e_date: <data>, 
    cust: { 
        name: <name>, 
        id: <id> 
    } 
}

rent:s_date = [ id_rent1, id_rent2, ... ]
\end{verbatim}

The first operation is instantaneous thanks to the index on \texttt{s\_date}.
The second operation is instantaneous thanks to \texttt{rent:id}.
The third operation is very efficient, as in the second case.
The fourth operation requires deleting the booking from both the first and second namespaces. Therefore, a lookup on \texttt{rent:id} and one on \texttt{rent:s\_date} are both needed, and both are fast.

\vspace{2mm}
\textbf{Option 3}

\begin{verbatim}
rent:s_date:id = { 
    type: <type>, 
    e_date: <date>, 
    cust: { 
        name: <name>, 
        id: <id> 
    } 
}
\end{verbatim}

The first operation is instantaneous.
The second operation is extremely inefficient: a linear scan must be performed.
The third operation is likewise inefficient.
The fourth operation is likewise inefficient.

\vspace{2mm}
\noindent
So, option 2 is the best choice.
\end{example}

\section{Document Databases}

Document databases store data as \keyword{documents}, which are self-contained, hierarchical data structures (typically JSON, BSON, or XML).

\begin{definition}{Document and Collection}
    A \keyword{document} is a single object where all attributes are stored together in a nested structure. A \keyword{collection} is a logical grouping of documents. Unlike relational tables, documents within the same collection are not required to have the same structure, although they should share some common characteristics.
\end{definition}

\subsection{Flexibility and Schema Design}

This structural flexibility is a major advantage over relational databases, allowing rapid iteration and easy accommodation of heterogeneous data. However, some discipline is beneficial:

\begin{itemize}
    \item Documents in a collection should share a \textbf{common base structure}.
    \item \textbf{Type attributes} are typically used to distinguish subtypes within a collection. For example, all \code{product} documents might share attributes like \code{name} and \code{price}, while a \code{type} field (\q{book}, \q{electronics}, \q{clothing}) indicates which additional, type-specific attributes are present.
    \item One can measure \keyword{collection cohesion} (how closely related the documents within a collection are) and \keyword{coupling between collections} (how dependent different collections are on each other) to decide whether to split a collection or merge two collections.
\end{itemize}

\subsection{Modeling Relationships}

Differently from relational databases, the One-to-Many relationship is more nuanced: it is split into three variants, based on the given cardinalities of the two parts~\cite{mongodbschemadesign}.

\begin{itemize}
    \item \textbf{One-to-Few Relationships:} The \q{one} is treated as the primary element, and the \q{many} side entities are \keyword{embedded} as an array within the primary document. For example, a \code{customer} document may embed an array of \code{addresses}. This allows retrieving the customer and all their addresses with a single read. However, one should ponderate the embedding choice based on whether the embedded element should have a separate document. This choice is typically driven by the most frequent queries on the system.
    \item \textbf{One-to-Many Relationships:} Each element of the \q{many} side gets referenced in a list inside the \q{one} side. Each element of the \q{many} side has now its own document. 
        For example, a \code{Product} may embed an array of \code{ReplacementPart} identifiers. In this way, to retrieve additional information about the the replacement parts of a product, more than just one queries will be needed. 
    \item \textbf{One-to-Squillions Relationships:} The \q{one} is related to so many \q{many} that the document is very likely to overflow the document size limit (16MB in MongoDB). So, we create a separate collection to contain the \q{many} documents, each of which will reference the \q{one}.
        For example, in a log-tracking system, a \code{host} produces thousands of \code{log}s per day. We create a collection for the \code{host}s and another one for the \code{log}s, where each \code{log} references its \code{host}.
    \item \textbf{Many-to-Many Relationships:} Two collections are used, and each collection contains \keyword{references} to instances of the other entity (Two-Way referencing). This minimizes data duplication, but referential integrity is not automatically guaranteed by the database --- it must be enforced at the application level. For example, \code{students} and \code{courses} collections each contain arrays of references to the other collection.
\end{itemize}

\subsection{Managing Hierarchies}

To represent hierarchical structures (e.g., a hierarchy of product categories in an e-commerce platform), multiple modeling approaches are available. The choice depends on whether read or write operations dominate, and on the hierarchy's depth and mutability.
    Consider an e-commerce product hierarchy where categories can have subcategories, forming a tree structure: \code{Electronics $\to$ Computers $\to$ Laptops}. The following approaches can model this:

    \begin{enumerate}
        \item \textbf{Parent References:} Each node stores a reference to its immediate parent. This makes moving a node easy (update one reference) but requires traversing the hierarchy step-by-step to retrieve all ancestors or descendants.
        \item \textbf{Children References:} Each node stores an array of references to its immediate children. This facilitates retrieving a node's direct children but makes finding all ancestors difficult.
        \item \textbf{Listing All Ancestors (Materialized Path):} Each node stores an array containing references to \emph{all} of its ancestors, from the root down to the immediate parent. This allows fast hierarchy navigation (retrieving all ancestors in a single read) and efficient subtree queries, at the cost of additional write overhead when the hierarchy is restructured.
    \end{enumerate}

\subsection{Denormalization}
We have talked about the different possibilities regarding relationships and hierarchies. Now, it's the time to explain the peculiar characteristic of document databases: \keyword{denormalization}.
Denormalization is the technique of adding redundancy to the data scheme, with the goal of making the queries more efficient. It is the inverse operation of normalization, which is a typical feature of relational databases.
Denormalization is driven by the cardinality of the entities and the frequency of the queries that the system must handle.
Let's consider an example.

\begin{example}{}
    An information system must manage data related to weather stations and their measurements.
    Each station makes millions of measurements over time, which makes it impossible to store a list of measurements in the document related to the station.
    Each station has a name, an installation city, an altitude, and an installation date (e.g., station \q{MI01} is located in Milan, at an altitude of 120 meters, and was installed in 2010).
    Each measurement has a timestamp, a temperature, and a humidity level.

    It is expected that the most frequent operations will be the following:
    \begin{description}
        \item[Q1]: Given a day, show all the measurements taken;
        \item[Q2]: Given the identifier of a station, show its last 100 measurements;
        \item[W1]: Insertion of a new measurement for a station.
    \end{description}

    Given the premises, we find ourselves in a One-to-Squillions relationship between a station and its measurements.
    Therefore, it is necessary to have two collections: Stations, Measurements.
    Having a separate collection for measurements allows query Q1 to be more efficient compared to a schema in which the measurements are embedded in the hosts that produce them.

    Query Q2 suggests a denormalization that would greatly improve its performance: we embed the last 100 measurements of a host within that host's document, in a list.
    With each new measurement, the oldest measurement will be discarded from this list and the new one will be added. The insertion of a new measurement (operation W1) therefore involves both the creation of a new document in the Measurements collection and the update of a list in a host document within the Hosts collection.

    Let us analyze the cost of the described denormalization: suppose there are 21 stations and that each of them produces 1 measurement per second. In this way, the daily frequency for operation W1 would be \(86400 \cdot 21 \approx 2000000\).
    Suppose also that queries Q1 and Q2 have a daily frequency of \(500000\) and \(1000000\), respectively.

    Let us calculate the cost of the three most frequent operations to understand whether denormalization is beneficial on a daily basis:
    \begin{description}
        \item[Q1]: It is sufficient to keep track of the first measurement of the day to obtain an efficient query. Let us roughly assume that the read cost is just one access.\footnote{This assumes that the Measurements collection follows an append-only policy, where documents are stored in insertion order. If a reference to the first measurement of each day is maintained, Q1 can locate the starting point with a single access and then stream the bulk of the day's measurements without scanning the entire collection. The cost of streaming the results is omitted here, as it is identical in both the denormalized and non-denormalized schemas and therefore does not affect the comparison.} 
            Therefore, \(2000000\) (two million) accesses per day.
        \item[Q2]: With the described denormalization, this query has the cost of a single access, therefore \(1000000\) (one million) accesses per day. Without denormalization, the cost is \(100\) accesses to the Measurements collection, which is a generous underestimate since it would be necessary to scan the Measurements collection. Therefore, the daily cost without denormalization is \(100000000\) (one hundred million) accesses.
        \item[W1]: Without denormalization, the cost of inserting a measurement is \(1\) write (we assume that the collection inserts documents in append mode), so the daily cost is \(20000000\) (twenty million) writes. With denormalization, to insert a measurement, a write to the Measurements collection, an access to the Stations collection, and a write to the station document are required. The daily cost is therefore \(2\cdot 20000000 = 40000000\) (forty million) writes and \(20000000\) (twenty million) accesses.
    \end{description}
    Therefore, without denormalization, we get \(2000000 + 100000000 = 102000000\) (one hundred two million) accesses and \(20000000\) (twenty million) writes. With denormalization, we get \(2000000 + 1000000 + 20000000 = 23000000\) (twenty-three million) accesses and \(40000000\) (forty million) writes.
    Let us not be fooled by the numbers: a write costs about as much as ten reads. Assuming this, without denormalization we get \(302000000\) (three hundred two million) reads, while with denormalization we get \(423000000\) (four hundred twenty-three) reads.
    According to this cost model, denormalization is not advantageous.
\end{example}

The example shows that making the choice of denormalization depends on the specific system and must be based upon deep analysis.
In general, denormalization is convenient of read-heavy systems and inconvenient on write-heavy systems.
But be careful: read- and write-heavy are not boolean, they are a continuous spectrum.

\section{Column Databases}

Column databases, also known as wide-column stores, organize data into \keyword{column families} --- groups of data items that are frequently accessed together. Columns within the same family are stored physically close to each other on disk, enabling efficient retrieval of related attributes without reading entire rows.

\subsection{An Example of Column Database: Bigtable}

Google's \keyword{Bigtable} is an influential column database. In Bigtable, a data value is indexed by the triple
\[
(\text{row identifier}, \text{column name}, \text{timestamp})
\]
The \keyword{timestamp} orders multiple versions of the same column value. When a new value is written, it is \emph{added} rather than overwriting the previous one, preserving a history of changes. Read and write operations are \keyword{atomic} even when they involve several columns within the same row. Rows are maintained in \keyword{sorted order} by their row identifier, which enables efficient range queries. However, the sort order can be specified in one way only --- there is no secondary ordering mechanism.

\subsection{Another Example: Cassandra}

\keyword{Apache Cassandra} is a wide-column database where columns can be added dynamically at any time, providing substantial schema flexibility.

\subsubsection{Architecture and Machine Hierarchy}

Cassandra's architecture follows a clear machine hierarchy:
\[
\text{Cluster} \to \text{Datacenter} \to \text{Rack} \to \text{Node}
\]
Data is distributed across nodes according to this hierarchy, with awareness of physical topology to ensure fault tolerance.

\subsubsection{Data Model}

\begin{itemize}
    \item \textbf{Column:} Composed of a \emph{name}, a \emph{value}, and a \emph{timestamp}. The timestamp is used to resolve conflicts between nodes when data is replicated.
    \item \textbf{Row:} An ordered collection of columns. A row must fit entirely within a single node.
\end{itemize}

\begin{definition}{Primary Key in Cassandra}
    In Cassandra, a \keyword{primary key} is not a single column but a list of columns. If the primary key consists of more than one column, it is called a \keyword{compound primary key}.
    \begin{itemize}
        \item The \keyword{partition key} is the first part of the primary key. It determines which node in the cluster stores the row.
        \item The \keyword{clustering key(s)} are the remaining columns of the primary key (after the partition key). They sort the data within a partition.
    \end{itemize}
\end{definition}

\begin{definition}{Table, Column Family, and Keyspace}
    A \keyword{table} is a container for rows and corresponds to a \keyword{column family}. If multiple column families exist, the row key contains a pointer to each one of them.

    A \keyword{keyspace} is an instance of a Cassandra database and is a container of tables. A keyspace is defined by:
    \begin{itemize}
        \item \textbf{Replication factor:} How many copies of each piece of data should exist.
        \item \textbf{Placement of replicas:} The strategy for distributing replicas across the cluster.
        \item \textbf{Set of tables:} The tables belonging to the keyspace.
    \end{itemize}
\end{definition}

\begin{warning}{Clock Synchronization}
    Cassandra relies on timestamps to resolve write conflicts between replicas, assuming a synchronized global clock across all nodes. Achieving this is \emph{not trivial} in distributed systems.

    Clock synchronization protocols such as \textbf{NTP} (Network Time Protocol) are typically employed. NTP allows nodes to synchronize their clocks with a hierarchy of reference time sources (stratum servers). However, even with NTP, clocks may drift between synchronization intervals, and network latency introduces uncertainty. For conflict resolution, Cassandra uses a \q{last-write-wins} policy based on these timestamps, which can lead to data loss if clock skew is significant. Alternative approaches, such as vector clocks or hybrid logical clocks, exist but introduce additional complexity.
\end{warning}

\subsubsection{Peer-to-Peer Communication and Anti-Entropy}

Cassandra uses a peer-to-peer architecture, where each peer communicates with up to three other peers using the \keyword{Gossip protocol}. This protocol disseminates cluster metadata (node state, topology changes) efficiently without a central coordinator --- here is why it is called \emph{gossip}.

To detect and repair data inconsistencies between replicas, Cassandra employs an \keyword{anti-entropy} mechanism:
\begin{enumerate}
    \item The sender and receiver nodes each calculate a \keyword{Merkle tree} over the data they store for a given range of tokens.
    \item The trees are compared. 
    \item If the trees do not match, the sender determines which node has the most recent version (using timestamps) and sends the up-to-date data to the receiver.
\end{enumerate}

\begin{note}{Merkle Trees}
    A \keyword{Merkle tree} (or hash tree) is a tree data structure in which each leaf node is labeled with the cryptographic hash of a data block, and each non-leaf node is labeled with the hash of the concatenation of its children's labels. The root of the tree --- the \keyword{Merkle root} --- serves as a compact, fixed-size digest of the entire dataset.

    \begin{center}
    \begin{tikzpicture}[
        level 1/.style={sibling distance=60mm},
        level 2/.style={sibling distance=30mm},
        every node/.style={draw, rectangle, minimum size=8mm, font=\footnotesize, rounded corners=2pt, fill=notecolor!10}
    ]
    \node {$H_{0,0} = H(H_{1,0} \parallel H_{1,1})$}
        child { node {$H_{1,0} = H(H_{2,0} \parallel H_{2,1})$}
            child { node {$H_{2,0} = H(D_0)$} }
            child { node {$H_{2,1} = H(D_1)$} }
        }
        child { node {$H_{1,1} = H(H_{2,2} \parallel H_{2,3})$}
            child { node {$H_{2,2} = H(D_2)$} }
            child { node {$H_{2,3} = H(D_3)$} }
        };
    \end{tikzpicture}
    \end{center}
    In Cassandra, the data on each node is partitioned into ranges. For each range, both the sender and receiver nodes build a Merkle tree whose leaves are hashes of individual rows. They compare the trees starting from the root:
    \begin{enumerate}
        \item If the roots match, the data for that range is consistent --- no further comparison is needed.
        \item If the roots differ, they compare the children. This proceeds recursively until the specific buckets (or rows) with discrepancies are identified.
    \end{enumerate}
    This allows detecting differences without transmitting or comparing entire datasets, drastically reducing network I/O during anti-entropy repair.

    Merkle trees are widely used in cryptographic systems to efficiently and securely verify the integrity and membership of data blocks. In Bitcoin and other cryptocurrencies, transactions within a block are organized into a Merkle tree. The Merkle root is stored in the block header, allowing lightweight clients to verify whether a transaction is included in a block without downloading the entire block --- this is called a Merkle proof.
\end{note}

While Cassandra has no central coordinator for cluster management, it does use a coordinator node for individual client requests (reads/writes). This is a temporary role, not a permanent or central one.
For each client request, one node acts as the coordinator --- usually the node the client connects to.
The coordinator:
\begin{enumerate}
    \item Routes the request to the appropriate replica nodes (based on the partition key).
    \item Collects responses from replicas and return the result to the client (error or success).
    \item Handles consistency (e.g., ensuring \code{QUORUM} responses are received).
\end{enumerate}

\keyword{Hinted handoff} is a complementary repair mechanism in Cassandra. When a node that should receive a write is temporarily unavailable, the coordinator stores a \q{hint} containing the write on behalf of the unavailable node. When the failed node comes back online, the hint is replayed, and the node catches up on missed writes. Hints are kept for a configurable time window; if the node remains down beyond that window, anti-entropy repair must be used instead.

\subsection{Guidelines for Column Databases}

Column databases should only be chosen when a large number of servers is required and the operational burden of managing a distributed column store is justified. Otherwise, one should opt for document databases or key-value stores, which offer similar data model flexibility without the same architectural complexity.

The following design guidelines should be observed when working with column databases:

\begin{itemize}
    \item \textbf{Design starts from queries.} Use cases guide the design of the data model. An analysis of queries yields information about: \emph{Entities} which become rows, \emph{Attributes of entities} which become columns, \emph{Query criteria} which drive the choice of column families and clustering keys, \emph{Derived values} which may be precomputed and stored.
    
    \item \textbf{Denormalize aggressively.} Since joins are not feasible in column databases, duplicate data across tables to support different query patterns.
    
    \item \textbf{Use valueless columns as flags.} A column with an empty value can act as a boolean flag, avoiding the need for explicit boolean values. One can freely mix valueless columns and columns with values within the same row.
    
    \item \textbf{Store an entire object in one row.} Related data should be colocated within the same row to enable atomic operations and single-pass retrieval.
    
    \item \textbf{Distribute load across nodes} by choosing an appropriate partition key that spreads data evenly across the cluster and avoids hot spots.
    
    \item \textbf{Set an appropriate limit on value versions.} Limiting the number of timestamped versions per column prevents unbounded storage growth.
    
    \item \textbf{Avoid complex column values.} BLOBs, JSON, and other complex structures should be split into individual columns instead, unless the application \emph{always} needs to retrieve the entire complex value as a single unit.
\end{itemize}

\section{Graph Databases}

Graph databases are designed to store and query entities and their relationships directly as nodes and edges. Several types of graph databases exist, each suited to specific use cases:

\begin{itemize}
    \item \textbf{Online vs. Offline graph databases:} Online graph databases support real-time, transactional querying and updates (OLTP, On-Line Transaction Processing), while offline graph databases are optimized for batch processing and analytical workloads (OLAP, On-Line Analytical Processing) over large, static graphs.
    \item \textbf{Native vs. Non-native graph databases:} A \keyword{native graph database} stores and manages data using a graph model at the storage layer (index-free adjacency). A \keyword{non-native graph database} uses an underlying non-graph storage engine (e.g., relational tables) and layers a graph API on top. A DBMS that exposes graph CRUD methods can be considered a graph DBMS regardless of its internal storage engine.
\end{itemize}

\subsection{Storage Model}

The storage model of a graph database can take several forms:

\begin{itemize}
    \item \textbf{Native graph storage:} The database engine is designed from the ground up to store nodes, relationships, and properties as first-class objects, with direct physical pointers between related entities.
    \item \textbf{Relational tables:} The graph structure is mapped to relational tables (e.g., tables for nodes and tables for edges with foreign keys). A graph query engine translates graph queries into SQL joins.
    \item \textbf{RDF triples:} Data is stored as subject-predicate-object triples in a triple store. This is common for semantic web and linked data applications.
\end{itemize}

\subsection{Data Model}

The predominant data model in modern graph databases is the \keyword{Labeled Property Graph} (LPG). It consists of the following core elements:

\begin{definition}{Labeled Property Graph}
A \keyword{Labeled Property Graph} is made up of:
\begin{itemize}
    \item \textbf{Nodes (vertices):} Represent entities. Each node can have one or more \emph{labels} that categorize it (e.g., \code{Person}, \code{Employee}).
    \item \textbf{Relationships (edges):} Represent connections between nodes. A relationship consists of:
    \begin{itemize}
        \item A \emph{start node} and an \emph{end node}.
        \item A \emph{name (type)} that describes the nature of the connection (e.g., \code{KNOWS}, \code{WORKS\_FOR}).
        \item A \emph{direction}, indicating the orientation of the relationship.
    \end{itemize}
    Relationship instances can also have properties (key-value pairs).
    \item \textbf{Properties:} Key-value pairs attached to both nodes and relationships, storing attributes (e.g., \code{name: "Alice"}, \code{since: 2020}).
    \item \textbf{Labels:} Tags assigned to nodes to group them into sets (e.g., \code{:Person}, \code{:Product}).
\end{itemize}
\end{definition}

Several structural characteristics distinguish the labeled property graph model:

\begin{itemize}
    \item \textbf{Multi-graph:} Multiple relationships of the same type, direction, and properties can exist between the same pair of nodes. Each relationship instance is a distinct entity, allowing several versions of the exact same edge to coexist.
    \item \textbf{Self-loops:} A relationship can connect a node to itself.
    \item \textbf{Identity:} Identity is determined by implementation-specific identifiers, not by properties. Two relationships can have identical data (type, properties, endpoints) and still be distinct entities distinguished only by their internal ID. Property values \emph{cannot} be references to vertices or edges; such connections can only be represented through explicit relationships.
\end{itemize}

\subsection{Queries}

Queries in a graph database are fundamentally based on \textbf{graph traversal and pattern matching}. Rather than joining tables through foreign key relationships, a graph query specifies a pattern of nodes and edges to match against the stored graph. The DBMS then locates paths through the graph that satisfy the pattern and returns the matched elements.

Graph databases are extremely efficient at graph traversals due to their storage and indexing strategies. 
A key technique is \textbf{index-free adjacency}: each node physically stores direct references (pointers) to its adjacent nodes and edges on disk. 
To traverse from one node to a neighboring node, the database follows an in-memory or on-disk pointer directly to the target node's storage location. 
This is an \(O(1)\) operation per hop, \textbf{independent of the total size of the dataset}. 
In contrast, a relational database would need to perform an index lookup (typically \(O(\log n)\) per join) for each relationship traversal, and the cost compounds with each additional level of depth, making deep traversals prohibitively expensive.

Node insertions and deletions in graph databases are also generally efficient. 
Inserting a new node involves allocating storage for the node's properties and establishing direct pointer references to its edges. 
Deleting a node requires removing the node and updating the adjacency pointers of all nodes that had edges to it. 
Because relationships are stored as physical pointers rather than matched through external indexes, these operations remain local to the affected nodes and their immediate neighbors. 
The cost is proportional to the node's degree (number of incident edges), not to the total graph size.

\subsection{Query Languages: Cypher}

\keyword{Cypher} is a family of declarative query languages for property graphs. While the full language is proprietary to Neo4j, the openCypher project provides an open subset that is also implemented by other vendors.
Other query languages are G-CORE and GQL (Graph Query Language), an ISO-standardized language for property graphs, heavily influenced by Cypher and other existing languages. GQL is the first standardized graph query language.

\subsection{Syntax and Semantics of a Query}

A Cypher query consists of a chain of clauses, evaluated in order:

\begin{enumerate}
    \item A \code{MATCH} clause followed by a graph pattern. Variants include:
    \begin{itemize}
        \item \code{OPTIONAL MATCH}: matches the pattern if it exists, but does not filter out rows if it does not (like an outer join). Unmatched elements are set to \code{null}.
        \item \code{MANDATORY MATCH}: forces the pattern to match; if it fails, the entire query returns no results (this is the default behavior, but can be made explicit).
    \end{itemize}
    \item (Optionally) a \code{WHERE} clause following a \code{MATCH} or \code{WITH}, providing filter conditions that restrict the matched patterns.
    \item (Optionally) a \code{WITH} clause followed by (possibly aliased) expressions and aggregate functions. It acts as a pipeline operator, passing only specified results to subsequent clauses and ending the current subquery. This is used for intermediate aggregation, projection, and pagination.
    \item A \code{RETURN} clause followed by (possibly aliased) expressions and aggregate functions. It occurs exactly once, at the end of the outermost query, and specifies what data to output.
    \item (Optionally) modifier sub-clauses: \code{ORDER BY} (sort results), \code{LIMIT} (restrict number of results), and \code{SKIP} (skip a number of results). These can follow a \code{WITH} or a \code{RETURN} clause.
    \item (Optionally) a \code{UNION} (or \code{UNION ALL}) clause to combine the results of two complete queries, each with its own \code{RETURN}. \code{UNION} removes duplicates; \code{UNION ALL} retains them.
\end{enumerate}

Queries are syntactically not nested, but \textbf{chained}: each clause consumes the output of the preceding one. The order of clauses directly affects query semantics. Implementations may evaluate clauses in a modified order as long as the semantics are preserved. For example, a query engine might reorder filter operations to reduce intermediate result sizes, provided the final result remains identical.

\begin{definition}{Node pattern}
    A \keyword{node pattern} is denoted by a pair of parentheses \code{()}. Optionally, the following components can be specified:
    \begin{itemize}
        \item A \textbf{variable name}, written directly inside the parentheses: \code{(foo)}.
        \item A \textbf{list of labels}, each prefixed with a colon \code{:} (e.g., \code{(:Person:Employee)}).
        \item A \textbf{set of properties}, written as a comma-separated list of key-value pairs inside curly braces \code{\{\}} (e.g., \code{\{name: 'Alice', age: 30\}}).
    \end{itemize}
\end{definition}

\noindent
Variable names, labels, and property keys are quoted with backticks (\code{`}), which can be omitted if only alphanumeric characters and underscores are used. Property values, if strings, use straight single quotes (\code{'}) or double quotes (\code{"}).

\begin{example}{Detailed Node Pattern}
    The following node pattern matches a node labeled \code{Person} and \code{Student}, with specific property constraints, bound to the variable \code{s}:
    \begin{center}
        \code{(s:\code{PERSON}:\code{STUDENT} \{name: 'Mario Rossi', age: 25, enrolled: true\})}
    \end{center}
\end{example}

\begin{definition}{Path Pattern and Relationship Pattern}
    A \keyword{path pattern} is a sequence of one or more node patterns separated by relationship expressions. A path can have three directions:
    \begin{itemize}
        \item \textbf{Forward:} \code{()-[...]->()}
        \item \textbf{Backward:} \code{()<-[...]-()}
        \item \textbf{Bidirectional} (undirected): \code{()-[...]-()}
    \end{itemize}

    \noindent
    The expression inside \code{[...]} is a \keyword{relationship pattern} with the following optional components:
    \begin{itemize}
        \item A \textbf{variable name}.
        \item A \textbf{list of relationship types}: a \code{|}-separated list of type names that begins with a colon \code{:} (e.g., \code{[:KNOWS|WORKS\_WITH]}).
        \item A \textbf{range literal} preceded by \code{*}: optionally followed by a specific range \code{n..m} (e.g., \code{*1..5}) or unbounded (e.g., \code{*}). The range specifies how many times the relationship may occur consecutively along the path. A specific range \code{*n..m} means the relationship type must occur between \code{n} and \code{m} times.
        \item A \textbf{set of properties}: a comma-separated list of \code{key:value} pairs inside curly braces \code{\{\}}.
    \end{itemize}
\end{definition}

\begin{warning}{}
    The range literal always applies to the whole relationship pattern. The \code{|}-separated list of relationship types is a \textbf{disjunction}: at least one of the listed types must match (in Neo4j, a relationship can only have one type). 
    In contrast, the comma-separated list of properties inside curly braces is a \textbf{conjuction}: the matched relationship must have \emph{all} of the specified properties.
\end{warning}

\begin{example}{Path Pattern with Range}
    The path pattern \code{(a)-[:FRIEND*1..3]->(b)} matches all paths that link two nodes with at least one and at most three consecutive \code{FRIEND} relationships: a path with three \code{FRIEND} relationships means that \code{(b)} is a friend of a friend of a friend of \code{(a)}.
\end{example}

Once nodes and relationships have been matched, their information can be accessed using the following functions and syntax:

\begin{itemize}
    \item Property access: \code{matched.propertyKey} or \code{matched[expression]}, where \code{expression} evaluates to a property key string;
    \item \code{id(matched)}: Returns the numeric internal identifier of a node or relationship. Identity is determined by this ID, not by properties;
    \item \code{keys(matched)}: Returns the list of all property keys of a node or relationship;
    \item \code{labels(node)}: Returns the list of labels assigned to a node;
    \item \code{type(relationship)}: Returns the type (name) of a relationship as a string.
\end{itemize}

\noindent
A matched path (alternating sequence of nodes and relationships) can be assigned to a variable: 
\begin{lstlisting}
MATCH path = <path_pattern> 
\end{lstlisting}
This is useful to extract the path informations with the following functions:

\begin{itemize}
    \item \code{relationships(path)}: Returns the list of relationships in the path.
    \item \code{nodes(path)}: Returns the list of nodes in the path.
    \item \code{length(path)}: Returns the number of relationships in the path.
\end{itemize}

\noindent
Since we mentioned lists, we now explain how to manipulate them.
To modify the elements of a list, list comprehension is used:
\begin{lstlisting}
<list_compr> ::= \[<item> IN <list> [WHERE <predicate>][| <map_expr>]\]
\end{lstlisting}
where \texttt{predicate} is a whatever boolean condition applied to the \texttt{item} and \texttt{map\_expr} is the expression to which the \texttt{item} will be mapped.

Elements of a list can be transformed to result rows by means of the \texttt{UNWIND} operator:   
\begin{lstlisting}
UNWIND <list | list_compr> AS <name>
\end{lstlisting}

\begin{example}{}
    For example, the following Cypher query returns all the friends of friends that connect two people.
    \begin{lstlisting}
MATCH path = (a)-[:FRIEND*]-(b)
UNWIND [node in nodes(path) | node.name] AS friend
RETURN friend
    \end{lstlisting}
\end{example}

\noindent
\begin{definition}{Graph pattern}
    A \keyword{graph pattern} is a comma-separated list of path patterns, interpreted in \textbf{conjunction}: all specified path patterns must be satisfied simultaneously for a match to occur.
\end{definition}

\noindent
In a graph pattern match:
\begin{itemize}
    \item Every node pattern is mapped to a node.
    \item Every path pattern is mapped to an alternating sequence of nodes and relationships (\texttt{NODE--RELATIONSHIP--NODE--\dots}), called a path.
    \item Matches of the same path pattern can include patterns of different lengths, and each distinct match of the pattern is called a \keyword{distinct match}.
    \item Parts of the pattern that have no associated variable must also be matched (they contribute to the matching conditions but are not bound to variables for later use).
\end{itemize}

\noindent
In a graph pattern match, each \textbf{relationship} (edge) must be used only once (\keyword{unique edges} feature). 
Each distinct path (without repeated relationships) counts as one result (\keyword{path multiplicity} feature).\footnote{This means that even if there are multiple ways to walk the same set of edges (by reordering intermediate nodes that are visited via different routes), a path is defined by its sequence of relationships. Two paths that use the \emph{same} relationships in the \emph{same} order are considered identical, yielding a single match. Two paths that use different relationships (or the same relationships in a different order) are distinct matches, each yielding a separate result.}

\begin{example}{}
    The following graph represents companies and which company sells to other ones.
    \begin{center}
    \begin{tikzpicture}[
        node distance=1.5cm,
        % Node styles
        concept/.style={circle, draw, minimum size=0.6cm, align=center, fill=red!10},
        property/.style={rectangle, draw, minimum size=0.8cm, align=center, fill=blue!10},
        % Edge styles
        relation/.style={-{Stealth[length=2mm]}, thick},
        propedge/.style={dotted, thick}
    ]
        % Nodes
        \node[concept] (a) at (0,0) {a};
        \node[concept, below left=1.5cm and 0.5cm of a] (b) {b};
        \node[concept, below right=1.5cm and 0.5cm of b] (c) {c};
        \node[concept, right=2.5cm of c] (d) {d};

        % Properties
        \node[property, right=0.2cm of a] (aprop) {name: 'Tubotec'};
        \node[property, right=0.2cm of b] (bprop) {name: 'Martelli'};
        \node[property, left=0.2cm of c] (cprop) {name: 'Burago'};
        \node[property, right=0.2cm of d] (dprop) {name: 'Pirolli'};

        % Relationships
        \draw[relation] (a) -- node[left] {:SELLS\_TO} (b);
        \draw[relation] (b) -- node[left] {:SELLS\_TO} (c);
        \draw[relation] (c) -- node[right] {:SELLS\_TO} (b);
        \draw[relation] (c) -- node[below] {:SELLS\_TO} (d);

        % Property connection
        \draw[propedge] (a) -- (aprop);
        \draw[propedge] (b) -- (bprop);
        \draw[propedge] (c) -- (cprop);
        \draw[propedge] (d) -- (dprop);

    \end{tikzpicture}
    \end{center}

    We want to find two companies which are connected with a path of four \texttt{:SELLS\_TO} relationships.
    The Cypher query is:
    \begin{lstlisting}
MATCH (x)-[:SELLS_TO*4]->(y)
RETURN x.name, y.name
    \end{lstlisting}
    However, it will not yield any result due to the unique edge feature.
    In fact, such a path of length 4 exists (\texttt{a->b->c->b->c}), but it requires the edge from \texttt{b} to \texttt{c} to be traversed twice.

    While if we want to find all such paths but of length 3, the related query would yield the following results:
    \begin{lstlisting}
x.name   |  y.name
--------- ----------
'Tubotec' 'Pirolli'
'Tubotec' 'Martelli'
    \end{lstlisting}
    for the distinct paths \texttt{a, b, c} and \texttt{a, b, d}, respectively.
\end{example}

\noindent
To use the same \textbf{relationship} multiple times, one must use several separate \code{MATCH} clauses, namely separate graph patterns.

\begin{example}{}
    Consider the graph of the previous example.
    If we want to reuse that edge from \texttt{b} to \texttt{c}, we can use add a separate \code{MATCH} clause:
    \begin{lstlisting}
MATCH (x)-[:SELLS_TO*3]->(y)
MATCH (y)-[:SELLS_TO]->(z)
RETURN x.name, y.name, z.name
    \end{lstlisting}
    The first \code{MATCH} clause will match with the distinct paths \texttt{a->b->c->b} and \texttt{a->b->c->d}.
    The second \code{MATCH} clause, for each of the previously matched paths, will reuse the bound \code{y} variable --- which is \code{b} for the first matched path and \code{d} for the second one --- and match companies that \code{y} sells to.
    For the first previous match, the path \texttt{b->c} will be matched --- we have legally reused the edge! --- while no match will be found for the second previous match since \code{d} has no outgoing edges.
    So, the Cypher query yields a single result:
    \begin{lstlisting}
 x.name  |  y.name  | z.name
--------- ---------- --------
'Tubotec' 'Martelli' 'Burago'
    \end{lstlisting}
\end{example}

\noindent
In a graph pattern match, each \textbf{node} can be used multiple times within the same graph pattern.

\begin{example}{}
    Consider the previous graph from the example.
    We want to find all the companies that Tubotec sells to via Burago. We can use two comma-separated path pattern to express this (but it can be expressed also with a single path pattern):
    \begin{lstlisting}
MATCH ({name: 'Tubotec'})-[:SELLS_TO*]->(x {name: 'Burago'}),
      (x)-[:SELLS_TO]->(y)
RETURN y.name
    \end{lstlisting}
    This query will yield two results:
    \begin{lstlisting}
  y.name
----------
'Martelli'
'Pirolli'
    \end{lstlisting}
\end{example}

\begin{note}{}
    Graph pattern matching in Cypher is based on \emph{relationship isomorphism}: relationships are matched injectively (no edge is reused within a single pattern), while nodes may be reused. This differs from pure homomorphism-based matching used in some other query languages (e.g., SPARQL), where both nodes and edges can be reused arbitrarily. The edge-isomorphism semantics prevent trivial, semantically redundant matches that would arise from re-walking the same edge.
\end{note}

\subsubsection{Path Quantifiers}

In addition to the range literal syntax for varying-length paths (e.g., \code{[:KNOWS*1..5]}), Cypher supports \keyword{path quantifiers}, also known as \keyword{quantified path patterns} or \keyword{quantified relationships}. These allow expressing complex, repeating sub-patterns along a path, going beyond simple repetition of a single relationship type.

\begin{definition}{Path Quantifier}
    A \keyword{path quantifier} applies a repetition operator to a parenthesized path pattern, specifying that the enclosed pattern may occur a certain number of times consecutively along a matched path. The syntax is:
    \[
    \texttt{((subpattern))\{min,max\}}
    \]
    where the quantifier \code{\{min,max\}} specifies the minimum and maximum number of repetitions of the subpattern. Variants include:
    \begin{itemize}
        \item \code{\{n\}}, exactly $n$ repetitions.
        \item \code{\{n,\}}, at least $n$ repetitions.
        \item \code{\{,m\}}, at most $m$ repetitions.
        \item \code{\{n,m\}}, between $n$ and $m$ repetitions.
        \item \code{\{*\}} or \code{*}, zero or more repetitions.
        \item \code{\{+\}} or \code{+}, one or more repetitions.
    \end{itemize}
\end{definition}

Path quantifiers are a generalization of the range literal in a relationship pattern. A simple path pattern \code{()-[:KNOWS*1..5]-()} is equivalent to the quantified path pattern \code{(()-[:KNOWS]-())\{1,5\}}. However, path quantifiers can enclose arbitrarily complex subpatterns involving multiple nodes and relationships, making them strictly more expressive than plain range literals.

Variables declared within a quantified path pattern are \keyword{group variables}. Rather than binding to a single node or relationship, they bind to lists of all the nodes or relationships they matched across the repetitions of the subpattern.

\begin{example}{}
    Path patterns can contain a \texttt{WHERE} clause. We can leverage this feature to show that path quantifiers are more expressive that range literals.
    Consider the following graph representing some of the places in Naples and their air distance:
\begin{center}
\begin{tikzpicture}[
    node distance=1.5cm,
    % Node styles
    concept/.style={circle, draw, minimum size=0.6cm, align=center, fill=red!10},
    property/.style={rectangle, draw, minimum size=0.8cm, align=center, fill=blue!10},
    % Edge styles
    relation/.style={-{Stealth[length=2mm]}, thick},
    propedge/.style={dotted, thick}
]
    % Nodes
    \node[concept] (a) at (0,0) {a};
    \node[concept, below right=0.3cm and 7.0cm of a] (b) {b};
    \node[concept, below right=1.5cm and 0.5cm of a] (c) {c};
    \node[concept, above right=0.2cm and 2.5cm of c] (d) {d};

    % Properties
    \node[property, above=0.2cm of a] (aprop) {name: 'Piazza Amedeo'};
    \node[property, above=0.2cm of b] (bprop) {name: 'Piazza del Plebiscito'};
    \node[property, below left=0.2cm and 0.5cm of c] (cprop) {name: 'Piazza San Pasquale'};
    \node[property, below=0.5cm of d] (dprop) {name: 'Piazza dei Martiri'};

    % Relationships
    \draw[relation] (a) -- node[above] {:NEXT} (b);
    \draw[relation] (a) -- node[left] {:NEXT} (c);
    \draw[relation] (c) -- node[below] {:NEXT} (d);
    \draw[relation] (d) -- node[below] {:NEXT} (b);

    % Property connection
    \draw[propedge] (a) -- (aprop);
    \draw[propedge] (b) -- (bprop);
    \draw[propedge] (c) -- (cprop);
    \draw[propedge] (d) -- (dprop);

\end{tikzpicture}
\end{center}

    We want to find a path from Piazza Amedeo to Piazza del Plebiscito such that there is less than a km between a place and another, so that we can stop at some bar for a fresh drink of water. We want to know the intermediate stops, of course.
    We can elegantly express this with the following Cypher query:
    \begin{lstlisting}
MATCH (x {name: 'Piazza Amedeo'}),
      (x) (()-[r:NEXT]->(s) WHERE r.distance <= 1) (y),
      (y {name: 'Piazza del Plebiscito'})
UNWIND s AS stop
RETURN stop.name
    \end{lstlisting}
    where we have used the \texttt{UNWIND} operator to transform each element of the \code{s} list in a row named \texttt{stop}.
    The output of the query is:
    \begin{lstlisting}
        stop
---------------------
'Piazza San Pasquale'
'Piazza dei Martiri'
    \end{lstlisting}
\end{example}

\subsubsection{The \texttt{WITH} Clause}
The \texttt{WITH} clause is used for intermediate projection, aggregation and aliasing. In contrast with the \texttt{WHERE} clause, whose role is to exclude rows from the result, the \texttt{WITH} clause manipulates data for subsequent queries.

\begin{example}{}
    Consider the following graph representing teams, their twinnings and their supporters:
    \begin{center}
    \begin{tikzpicture}[
        node distance=1.5cm,
        % Node styles
        concept/.style={circle, draw, minimum size=0.6cm, align=center, fill=red!10},
        property/.style={rectangle, draw, minimum size=0.8cm, align=center, fill=blue!10},
        % Edge styles
        relation/.style={-{Stealth[length=2mm]}, thick},
        propedge/.style={dotted, thick}
    ]
        % Nodes
        \node[concept] (t1) at (0,0) {t1};
        \node[concept, right=4cm of t1] (t2) {t2};
        \node[concept, below right=3cm and 1.8cm of a] (t3) {t3};

        \node[concept, below left=2cm and 2cm of t1] (s1) {s1};
        \node[concept, left=3cm of t1] (s2) {s2};
        \node[concept, above=2cm of t1] (s3) {s3};

        \node[concept, above right=1cm and 1cm of t2] (s4) {s4};

        % Relationships
        \draw[relation] (t1) -- node[below] {:TWINNED} (t2);

        \draw[relation] (s1) -- node[right] {:SUPPORTS} (t1);
        \draw[relation] (s2) -- node[above] {:SUPPORTS} (t1);
        \draw[relation] (s3) -- node[right] {:SUPPORTS} (t1);

        \draw[relation] (s4) -- node[right] {:SUPPORTS} (t2);
    \end{tikzpicture}
    \end{center}
    We want to find the names of the ones that support a team with at least two supporters.
    \begin{lstlisting}
MATCH (t)<-[:SUPPORTS]-(s)
WITH t, count(s) AS supporters
WHERE supporters >= 2
MATCH (s)-[:SUPPORTS]->(t)
RETURN s.name AS supporter, t.name AS team
    \end{lstlisting}
\end{example}

\noindent
In this clause, we can also apply reductions --- \code{count}, \code{sum}, \code{avg}, \code{min}, \code{max}, \code{stDev}, \code{percentileDisc}, etc. ---, aliasing via the \texttt{AS} binary operator (as shown in the previous example), and a special form of aggregation with the \texttt{collect} function, that gathers rows into a list.

\begin{note}{Implicit Grouping}
    When a \code{WITH} or \code{RETURN} clause contains both aggregate and non-aggregate expressions, the non-aggregate expressions serve as implicit grouping keys.
    For example, in the previous example, at the point of the \texttt{WITH} clause, rows are grouped by the team node \texttt{t}.
\end{note}

\subsubsection{The \texttt{WHERE} Clause}

The \texttt{WHERE} clause determines which matched patterns are retained. A \code{WHERE} clause contains a boolean expression, which can be arbitrarily complex, i.e. made by connecting and nesting several \emph{basic expressions}. 
Basic expressions are: 
\begin{itemize}
    \item checks over properties and variables of the form \texttt{<var|property> <binary\_operator> <expr>} or \texttt{<var|property> <unary\_operator>}. The binary operators \texttt{=, <>} can be used with all kinds of variables and properties.
        \begin{example}{Checking Identity}
            Find two distinct people s.t. the first knows a computer scientist and the second is married with them.
            \begin{lstlisting}
MATCH (a)-[:KNOWS]-(x{occupation: 'Computer Scientist'}),
      (b)-[:SPOUSE]-(x)
WHERE a <> b
RETURN a, b, x
            \end{lstlisting}
            In Cypher, \texttt{a <> b} checks the node identity. To check the property equality, one has to compare all the nodes properties.
        \end{example}
        The unary operators \texttt{IS NULL} and \texttt{IS NOT NULL} are available for all kinds of variables and properties too.
        For strings, \texttt{STARTS WITH}, \texttt{ENDS WITH}, \texttt{CONTAINS} can be used, as well as to \texttt{=\~} match against regular expressions.
        For lists, there are specific functions that apply to list comprehension: 
        \begin{itemize}
            \item The function \texttt{all} returns true when the list comprehension predicate is true for \emph{all} the elements of the \texttt{list}. For example, \texttt{all(person IN nodes(path) WHERE person.age > 18)} evaluates to \texttt{true} if \texttt{n.age > 18} is true for every item.
            \item The function \texttt{any} returns true when the predicate is true for \emph{any} of the elements of the \texttt{list}.
            \item The function \texttt{none} returns true when the predicate is true for \emph{none} of the elements of the \texttt{list}.
            \item The function \texttt{single} returns true when the predicate is true only for a \emph{single} element of the \texttt{list}.
        \end{itemize}
    \item path patterns.
        \begin{example}{}
            Find the people who know a computer scientist and who are not married.
            \begin{lstlisting}
MATCH (a)-[:KNOWS]-({occupation: 'Computer Scientist'}),
WHERE NOT (a)-[:SPOUSE]-()
RETURN a
            \end{lstlisting}
            Without the \texttt{WHERE} clause, the query would match also married people.
        \end{example}
\end{itemize}
Expressions can be connected with binary logical operators --- \texttt{AND}, \texttt{OR}, \texttt{XOR} --- or negated with the unary operator \texttt{NOT}.

\subsubsection{Other Features}

Cypher provides several additional features for advanced querying:

\begin{itemize}
    \item Shortest path functions: \code{shortestPath((a)-[*]-(b))} returns a single shortest path between two nodes. \code{allShortestPaths((a)-[*]-(b))} returns all shortest paths.
    \item \code{DISTINCT}: It eliminates duplicate results. When applied in a \texttt{RETURN} or a \texttt{WITH} clause, it removes duplicate rows from the result set, and it can be used as well in aggregate functions (e.g., \texttt{COUNT(DISTINCT x)}.
\end{itemize}

Now, let's see two examples that comprise most of the introduced features of Cypher.
\begin{example}{Friend Recommendation}
    Consider a graph representing people and friend relationships. The following graph is a subgraph that shows an example of two people who have three mutual friends but who are not friends:
    \begin{center}
    \begin{tikzpicture}[
        node distance=1.5cm,
        % Node styles
        concept/.style={circle, draw, minimum size=0.6cm, align=center, fill=red!10},
        property/.style={rectangle, draw, minimum size=0.8cm, align=center, fill=blue!10},
        % Edge styles
        relation/.style={-{Stealth[length=2mm]}, thick},
        propedge/.style={dotted, thick}
    ]
        % Nodes
        \node[concept] (a) at (0,0) {a};
        \node[concept, right=6cm of a] (b) {b};
        \node[concept, below right=1.5cm and 3.0cm of a] (x) {x};
        \node[concept, right=2.8cm of a] (y) {y};
        \node[concept, above right=1.5cm and 3.0cm of a] (z) {z};

        % Relationships
        \draw[relation] (a) -- node[left] {:KNOWS} (x);
        \draw[relation] (a) -- node[above] {:KNOWS} (y);
        \draw[relation] (a) -- node[left] {:KNOWS} (z);
        \draw[relation] (x) -- node[above] {} (a);
        \draw[relation] (y) -- node[above] {} (a);
        \draw[relation] (z) -- node[above] {} (a);

        \draw[relation] (b) -- node[right] {:KNOWS} (x);
        \draw[relation] (b) -- node[above] {:KNOWS} (y);
        \draw[relation] (b) -- node[right] {:KNOWS} (z);
        \draw[relation] (x) -- node[above] {} (b);
        \draw[relation] (y) -- node[above] {} (b);
        \draw[relation] (z) -- node[above] {} (b);
    \end{tikzpicture}
    \end{center}
    We want to find all pairs of people who are not yet friends but share at least three common friends, and rank them by the number of common friends. Return up to 10 results, skipping the first 5.
    \begin{lstlisting}[language=SQL, caption={Friend recommendation query}]
MATCH (a:Person)-[:KNOWS]->(common:Person)<-[:KNOWS]-(b:Person)
WHERE NOT (a)-[:KNOWS]-(b)
WITH a, b, count(common) AS commonFriends
WHERE commonFriends >= 3
RETURN a.name, b.name, commonFriends
ORDER BY commonFriends DESC
SKIP 5
LIMIT 10
    \end{lstlisting}
\end{example}

\begin{example}{Path Existence and Aggregation}
    For each person, find, if possible, the shortest chain of \code{KNOWS} relationships to any person named \code{'Jessica'} who works at \code{'Tubotec'}. Return the person's name, the chain length, and the names along the path.
    \begin{lstlisting}[language=SQL, caption={Shortest path query with aggregation}]
MATCH (b:Person {name: 'Jessica'})-[:WORKS_AT]->(:Company {name: 'Tubotec'})
MATCH (a:Person)
OPTIONAL MATCH p = shortestPath((a)-[:KNOWS*]-(b))
WITH a, p, length(p) AS distance
ORDER BY distance ASC
WITH a, collect(p)[0] AS shortestPath, 
     collect(distance)[0] AS minDistance
RETURN a.name AS person,
       minDistance AS hopsToAlice,
       [n IN nodes(shortestPath) | n.name] AS pathNames
ORDER BY minDistance ASC
    \end{lstlisting}
    This query: 
    \begin{enumerate}
        \item matches, for each person \texttt{a}, the shortest path \texttt{p}. Rows are implicitly grouped by \texttt{a};
        \item projects \texttt{a} and \texttt{p} and the computed \texttt{distance};
        \item sorts rows by distance in ascending order;
        \item keeps \texttt{a}, the first path, the first distance;
        \item returns \texttt{a.name}, the distance, a list containing the names of the people in the path;
        \item sorts by distance in ascending order, so that the people who know Jessica directly will appear in the first rows.
    \end{enumerate}
\end{example}

\subsubsection{Error Handling}

Cypher uses the special value \code{null} to indicate missing, unknown, or erroneous values. The semantics of \code{null} in Cypher follow a three-valued logic:

\begin{itemize}
    \item Any operation or function that receives \code{null} as input typically returns \code{null} (e.g., \code{null + 5 = null}, \code{null = null \(\to\) null}).
    \item In boolean contexts, \code{null} is treated as false (falsy semantics): if a \code{WHERE} condition evaluates to \code{null}, the row is filtered out.
    \item A \code{null} value cannot be distinguished from an unbound variable: accessing a property that does not exist returns \code{null}.
\end{itemize}

\noindent
A Cypher query fails (raises an error) in the following situations:

\begin{itemize}
    \item Syntax errors: The query string is not valid Cypher syntax.
    \item Type errors: Certain operations require specific types and fail if given incompatible arguments (e.g., using a string where a number is required in a numeric function).
    \item Constraint violations: Attempting to create a node or relationship that violates a uniqueness constraint or other schema constraint.
    \item Resource exhaustion: Queries that consume excessive memory or time may be terminated by the database engine.
    \item Division by zero: In some implementations, an explicit division by zero raises an error.
\end{itemize}

When a query fails, it returns an error to the client rather than a partial result. A \code{null} in the result, by contrast, does \emph{not} indicate query failure --- it simply signals the absence of a value for that particular row or expression.

\subsection{Graph Database Design}

Designing a graph database requires a shift in mindset from relational modeling: instead of mapping entities to tables, attributes to columns and relationships to columns or tables, one must reason about nodes, labels, relationships and properties.

\subsubsection{Design Process}

The design process follows a structured sequence of steps\cite{robinson2015graph}:

\begin{enumerate}
    \item \textbf{Define Nodes.} Nodes map to the \textbf{domain objects} (entities) of the application. These are the \q{nouns} in the use cases: people, products, orders, accounts, devices, locations, etc. Each node represents a distinct instance of an entity.

    \item \textbf{Define Labels.} Labels are used to \textbf{group nodes into sets}, analogous to classes or types. A node can carry multiple labels, allowing flexible categorization. For example, a node representing a person who is both a customer and an employee can be labeled \code{:Person:Customer:Employee}. Labels enable efficient filtering during queries and can be indexed by the database engine.

    \item \textbf{Define Relationships.} Relationships model the \textbf{interactions and connections} between nodes. These correspond to the \q{verbs} in the use cases: \emph{knows}, \emph{purchased}, \emph{works\_for}, \emph{located\_in}. Each relationship must have:
    \begin{itemize}
        \item A \textbf{type} (name) that describes its semantics.
        \item A \textbf{direction} (though queries can traverse relationships in either direction).
    \end{itemize}
    Relationships should be fine-grained and semantically meaningful. Prefer specific types (\code{:PLACED\_ORDER}, \code{:CANCELLED\_ORDER}) over generic ones (\code{:HAS\_INTERACTION}) to make queries self-documenting and efficient.

    \item \textbf{Define Properties.} Properties are determined by the \textbf{expected queries} and the functional requirements of the database. Both nodes and relationships can carry properties. The guiding questions are:
    \begin{itemize}
        \item What attributes are needed to filter, sort, or return data in queries?
        \item What attributes are required by the business logic?
        \item What attributes must be indexed for performance?
    \end{itemize}
    Properties that are only occasionally accessed can be separated from the core entity, but denormalization is generally preferred over joins in graph databases.
\end{enumerate}

\noindent
Graph database design is inherently \emph{iterative}. 
One should start by modeling the entities and relationships that directly support the most critical queries, then test the model with realistic datasets and query workloads, then refine, adjusting label granularity, or denormalizing data to eliminate performance bottlenecks. 
The flexibility of the graph model makes incremental evolution far less costly than in relational systems.

\subsubsection{Properties vs. Relationships}

A key design decision is whether to model a data element as a \textbf{property} or as a \textbf{relationship} to a separate node.

\begin{example}{Property vs. Relationship}
    Consider modeling a person's city of residence. One could use:
    \begin{itemize}
        \item A \textbf{property}: \code{(:Person \{city: 'Naples'\})}.
        \item A \textbf{relationship} to a dedicated node: 
        \begin{center}
            \code{(:Person)-[:LIVES\_IN]->(:City \{name: 'Naples'\})}.
        \end{center}
    \end{itemize}
\end{example}
This is actually a normalization vs. denormalization problem: storing an information as a node property (denormalization) is better for reads --- the property can be directly accessed ---, but inconvenient for updates --- the property must be updated for each node, possibly causing inconsistency.
On the other hand, storing an information as a relationship is convenient for updates --- the node that contains the information gets updated once and consistency is guaranteed ---, while being inconvenient for reads --- one hop is added per node.

In general, relationships yield cleaner queries and better performance when:

\begin{itemize}
    \item The value is shared across many entities (avoid data duplication and enable efficient \q{who else lives in Naples?} queries).
    \item The value itself has attributes (a city has a population, region, country---a property string cannot capture this).
    \item The value is a filtering or aggregation criterion in queries (nodes can be indexed; property string matching is slower).
    \item The value participates in further traversals (\q{find all people living in cities connected by a direct train}).
\end{itemize}

\noindent
A property should remain a property when:

\begin{itemize}
    \item It is a simple, scalar attribute with no independent identity (e.g., age, birthdate, email).
    \item It is always accessed together with the owning node.
    \item There is no need to traverse from the value to other entities.
\end{itemize}

